\section{Conclusion}

This article presented the definitions and properties of Cycles, a fundamental
concept that enables representation of repeating sequences of values. We
defined Cycles using two approaches---a recursive definition and a modulo-based
definition---and proved their equivalence for all positions. We further
established eleven properties: element access via modular indexing, direct
access for small positions, invariance under addition of cycle-period
multiples, consistency across distinct multiples, modulo propagation from
values to cycle access, repeated-cycle invariance, value positivity, rotation
invariance, and the all-zero, none-zero, and some-zero transfers of
\texttt{MemCycle}'s divisor-residue classification from the base list to cycle
positions.

\begin{equation*}
\begin{aligned}
&\forall\,L\in\mathbb{L},\quad \forall\,v\in\mathbb{N}_0,
  \quad \forall\,i,m_1,m_2\in\mathbb{N}_0 \\
L&:=[v_0,v_1,\dots,v_{n-1}]\in\mathbb{N}_0^n,\quad |L|>0 \\
\operatorname{Cycle}&:=[v_0,v_1,\dots,v_{n-1},v_0,v_1,\dots] \\
n&=|L|
\end{aligned}
\end{equation*}

\begin{equation*}
\operatorname{RecCycle}_i=\operatorname{ModCycle}_i
  =\operatorname{MemCycle}_i
  \quad \text{[Three-Way Equivalence]}
\end{equation*}

\begin{equation*}
\begin{gathered}
\operatorname{Cycle}_i=L[i\operatorname{mod}n]
  \quad \text{[Cycle Element Access]} \\
\text{key}<n\implies
  \operatorname{Cycle}_{\text{key}}=L_{\text{key}}
  \quad \text{[Small Value Direct Lookup]} \\
\operatorname{Cycle}_i\operatorname{mod}d
  =\operatorname{Cycle}_{(i\operatorname{mod}n)}\operatorname{mod}d
  \quad \text{[Mod Propagation]}
\end{gathered}
\end{equation*}

\begin{equation*}
\begin{gathered}
\operatorname{Cycle}_{(i+n\cdot m)}=L[i\operatorname{mod}n]
  \quad \text{[Value Match After Many Loops]} \\
\operatorname{Cycle}_{(i+n\cdot m_1)}
  =\operatorname{Cycle}_{(i+n\cdot m_2)}
  \quad \text{[Two Multiples]} \\
\operatorname{Cycle}^{(t)}_{\text{pos}}
  =\operatorname{Cycle}_{\text{pos}}\quad \forall t>0
  \quad \text{[Repeated-Cycle Invariance]}
\end{gathered}
\end{equation*}

\begin{equation*}
\begin{gathered}
(\forall x\in L,\ x\geq 0)
  \implies \operatorname{Cycle}_{\text{pos}}\geq 0
  \quad \text{[Cycle Value Positivity]} \\
(\forall y\in L,\ y>x)
  \implies \operatorname{Cycle}_{\text{pos}}>x
  \quad \text{[Cycle Value Greater Than]} \\
\operatorname{rotateAt}(\operatorname{Cycle},k)_i
  =\operatorname{Cycle}_{i+k}
  \quad \text{[Rotation Invariance]}
\end{gathered}
\end{equation*}

\begin{equation*}
\begin{gathered}
\operatorname{allZero}(d)\implies
  \forall k,\ \operatorname{Cycle}_k\operatorname{mod}d=0
  \quad \text{[All-Zero Residue Transfer]} \\
\operatorname{noneZero}(d)\implies
  \forall k,\ \operatorname{Cycle}_k\operatorname{mod}d\neq 0
  \quad \text{[None-Zero Residue Transfer]} \\
\operatorname{someZero}(d)\implies \exists\,k_0,k_1,\quad
  \operatorname{Cycle}_{k_0}\operatorname{mod}d=0\ \land{} \\
\operatorname{Cycle}_{k_1}\operatorname{mod}d\neq 0
  \quad \text{[Some-Zero Residue Transfer]}
\end{gathered}
\end{equation*}

The reusable equivalence, access, periodicity, modulo, repetition, positivity,
and rotation properties were formally verified using Scala Stainless. The
residue-classification counts are likewise verified on the base list; the
all-zero, none-zero, and some-zero transfers stated here are the mathematical
corollaries proved in Subsections~5.10--5.12. The full Stainless verification
code is in Appendix~A.

\section{Future Work}

Future work may include exploring more complex properties of Cycles, such as
their behavior under operations including concatenation and filtering, and
their applications in algorithms and data structures. Additionally, we can
investigate discrete integration of Cycles, similar to the work done for
lists~\cite{mata2026lists} and integrals~\cite{mata2026integral}.
